% =============================================================================
% experiments.tex  --  drop-in replacement for \section{Experiments}
%
% All numbers are PLACEHOLDERS ("--"); fill from the result logs.
% Tables use only core LaTeX (\multicolumn, \hline) + \resizebox (graphicx).
%
% Structure: Main Results (recognition transfer + SPT vs B0 following/success,
%   sim + real) and Ablation Studies (the 2x2 over SPT's two components).
%   2x2 configs:  B0 = neither | PL = Stage-B conditioning only |
%   VQA = Stage-A prediction loss only | SPT = both (ours).
% =============================================================================

\section{Experiments}
\label{sec:experiments}

We evaluate SPT in simulation and on a real robot, and ask two questions.
(i)~Does the preference prediction learned on the source transfer to the target,
that is, can the frozen source model label target demonstrations correctly?
(ii)~Does Stage~B turn those labels into a policy that follows the commanded
preference on the target?

\subsection{Experimental Setup}
\label{sec:exp-setup}

\textbf{Simulation benchmark.}
We synthesize a preference-varied manipulation dataset in
RoboTwin~2.0~\citep{robotwin}, spanning five preference axes: \emph{grasp contact
height} (low vs.\ high on the object), \emph{release height} (low vs.\ high above
the target), \emph{obstacle clearance} (narrow vs.\ wide detour), \emph{grasp
orientation} (side vs.\ top), and \emph{placement location} (center vs.\ corner).
For each axis the source has eight tasks, each demonstrated under both preference
values with 100 demonstrations per value. The target is a held-out task with new
objects on the same axis; its instructions do not contain the preference, so
each target demonstration's preference is unknown at training time. Observations
are three Intel RealSense~L515 RGB views.

\textbf{Real-robot validation.}
We test on a UR5e under two cross-task settings. For \emph{obstacle clearance},
we label a near vs.\ far clearance preference on pick-and-place and transfer it
to pushing (push a bowl around an obstacle). For \emph{strike strategy}, we label
a direct vs.\ bank-shot preference on air hockey and transfer it to billiards
(cue the ball into a pocket directly or off a cushion). As in simulation, the
target instructions do not state the preference. Each target shares the source
axis but differs in objects, dynamics, and goal.

\textbf{Metrics.}
We report three numbers. \emph{Recognition accuracy} is how often the frozen
source predictor assigns the correct preference to a target demonstration
(chance $50\%$). \emph{Preference-following accuracy} is the fraction of rollouts
whose behavior matches the commanded preference, scored by a per-axis statistic
$\phi$ (contact height; release height; minimum obstacle clearance; wrist
orientation; placement location), with continuous axes thresholded from source
statistics and categorical axes matched directly. \emph{Task success} is the
completion rate regardless of preference. We fix all criteria and thresholds
before evaluation; details are in the appendix.

\textbf{Configurations.}
SPT has two components: the Stage-A prediction loss and the Stage-B conditioning
on the self-labels. We study both with a $2\times2$ over these factors.
\textbf{B0} uses neither and finetunes on the target with the action loss only.
\textbf{PL} conditions Stage~B on the self-labels but starts from a policy
trained without the prediction loss. \textbf{VQA} starts from the
prediction-trained model but does not condition Stage~B on the labels.
\textbf{SPT} uses both. The four share the backbone, data, and compute budget.

\textbf{Implementation details.}
The backbone is Qwen3-VL-4B. The action expert is an action-token head in the
OpenVLA-OFT style, trained with an $L_1$ loss on $14$-D joint action chunks of
horizon $50$, and warm-started from a manipulation-pretrained checkpoint;
training from scratch does not converge on this data. The prediction head is the
VLM's language head, supervised with single-token cross-entropy over a short
head-camera clip. We encode the active end-effector pose at the clip frames with
a small MLP $\phi$ ($9\!\to\!512\!\to\!h$) and inject it as a token at an
existing special-token slot, so the pretrained checkpoint loads without resizing
the embedding. Stage~A optimizes $\mathcal{L}_A$; Stage~B continues training the
action expert at a lower learning rate. The appendix lists all hyperparameters
and the $\phi$ and $w(\cdot)$ vocabularies.

\subsection{Main Results}
\label{sec:exp-main}

\textbf{Preference prediction transfers to the target.}
Table~\ref{tab:sim_recognition} gives the frozen predictor's accuracy on the
target, per axis. It is well above chance on all five axes, so the preferences
learned on the source carry over to the target's new tasks and objects, with no
preference stated in the instruction. These predictions are the labels Stage~B
trains on.

\begin{table}[t]
  \centering
  \small
  \caption{\textbf{Prediction transfer (simulation).} Accuracy of the frozen
  source predictor on the target, per preference axis (chance $=50\%$). It
  depends only on the Stage-A predictor, so it is reported once per axis.}
  \label{tab:sim_recognition}
  \begin{tabular}{lc}
    \hline
    Preference axis & Recognition Acc.\ (\%) \\
    \hline
    Grasp contact height & -- \\
    Release height       & -- \\
    Obstacle clearance   & -- \\
    Grasp orientation    & -- \\
    Placement location   & -- \\
    \hline
    Average              & \textbf{--} \\
    \hline
  \end{tabular}
\end{table}

\textbf{SPT follows the commanded preference.}
Table~\ref{tab:sim_main} reports preference-following and success in simulation
for SPT and the no-steering baseline B0. SPT follows the commanded preference on
every axis, while B0 keeps to one behavior; success is unchanged, so the added
control does not cost task completion. The real-robot results
(Table~\ref{tab:real_main}) match: SPT follows the preference on both targets
while B0 does not.

\begin{table}[t]
  \centering
  \small
  \caption{\textbf{Preference-following and task success (simulation).} Per-axis
  results for SPT and the no-steering baseline B0.}
  \label{tab:sim_main}
  \begin{tabular}{lcccc}
    \hline
    & \multicolumn{2}{c}{Following Acc.\ (\%)} & \multicolumn{2}{c}{Success (\%)} \\
    Preference axis & B0 & SPT & B0 & SPT \\
    \hline
    Grasp contact height & -- & -- & -- & -- \\
    Release height       & -- & -- & -- & -- \\
    Obstacle clearance   & -- & -- & -- & -- \\
    Grasp orientation    & -- & -- & -- & -- \\
    Placement location   & -- & -- & -- & -- \\
    \hline
    Average              & -- & \textbf{--} & -- & -- \\
    \hline
  \end{tabular}
\end{table}

\begin{table}[t]
  \centering
  \small
  \caption{\textbf{Real-robot results (UR5e).} Preference-following accuracy and
  task success on two cross-task settings, SPT vs.\ B0. $N$ is the number of
  rollouts per condition.}
  \label{tab:real_main}
  \begin{tabular}{llccccc}
    \hline
    & & & \multicolumn{2}{c}{Following Acc.\ (\%)} & \multicolumn{2}{c}{Success (\%)} \\
    Setting & Preference & $N$ & B0 & SPT & B0 & SPT \\
    \hline
    Obstacle clearance & near   & -- & -- & -- & -- & -- \\
    (push)             & far    & -- & -- & -- & -- & -- \\
    Strike strategy    & direct & -- & -- & -- & -- & -- \\
    (billiards)        & bank   & -- & -- & -- & -- & -- \\
    \hline
  \end{tabular}
\end{table}

\subsection{Ablation Studies}
\label{sec:exp-ablation}

We ablate SPT's two components with the configurations of
Section~\ref{sec:exp-setup}, averaged over the five axes
(Table~\ref{tab:ablation}). Conditioning Stage~B on the self-labels is what makes
the policy follow the preference: without it (VQA) the policy stays at the B0
level, and with it (SPT) it follows. The Stage-A prediction loss adds a further
gain: given the same self-labels, starting from the prediction-trained model
(SPT) follows better than starting from one trained without it (PL). Either
component alone is below SPT.

\begin{table}[t]
  \centering
  \small
  \caption{\textbf{Ablation of SPT's two components} (simulation, averaged over
  the five axes). \textbf{PL}: Stage-B self-label conditioning only;
  \textbf{VQA}: Stage-A prediction loss only; \textbf{SPT}\,=\,both
  (Section~\ref{sec:exp-setup}).}
  \label{tab:ablation}
  \begin{tabular}{lcc}
    \hline
    Configuration & Following Acc.\ (\%) & Success (\%) \\
    \hline
    B0 \,(neither)              & -- & -- \\
    PL \,(conditioning only)    & -- & -- \\
    VQA (prediction only)       & -- & -- \\
    SPT (both)                  & \textbf{--} & -- \\
    \hline
  \end{tabular}
\end{table}

\subsection{Limitations and Future Work}
\label{sec:limitations}

SPT inherits the source preference set. A user can only steer the target with
values $c\in\mathcal{C}$ that were defined and labeled on the source, and the
predictor only distinguishes those values. A new value on an existing axis, or a
new axis, is out of scope. Extending the set at test time, composing labeled
axes, and discovering preference axes from data are left to future work.
